\subsection{Zvolené prompty}
\label{subsec:choice-prompts}

Prompty jsou v souladu se strukturou popsanou v \secref[sekci]{subsec:prompts-structure} sestaveny z role modelu, pravidel pro výběr nálezů a pokynů pro formát výstupu.
Jsou psány v angličtině, protože jazykové modely dosahují konzistentněji kvalitních výsledků s anglickými instrukcemi.

\subsubsection*{One-shot prompt}
\label{subsubsec:choice-prompts-oneshot}

Při one-shot přístupu je analýza provedena jediným voláním modelu.
Systémový prompt (\lstref{lst:prompt-oneshot}) definuje roli modelu, typy hledaných problémů a požadovaný formát odpovědi -- kombinuje analytickou část z \lstref{lst:prompt-cot-draft} s formátovacími pokyny z \lstref{lst:prompt-cot-review}.

\subsubsection*{Chain-of-thought prompty}
\label{subsubsec:choice-prompts-cot}

Při chain-of-thought přístupu jsou volání modelu tři.
Každý krok má vlastní systémový prompt, který modelu zadává přesně ohraničenou úlohu.

Každý krok má vlastní systémový prompt, který modelu zadává přesně ohraničenou úlohu: hrubá analýza (\lstref{lst:prompt-cot-draft}), kritická revize (\lstref{lst:prompt-cot-critique}) a finální přehled (\lstref{lst:prompt-cot-review}).

Úplné znění všech promptů je uvedeno v příloze~\ref{ch:appendix-prompts}.

\endinput